% =====================================================================
% §1.6 — Integración de la deshidratadora en el flujo post-cosecha
% =====================================================================
\section{Integración de la deshidratadora en el flujo post-cosecha}
\label{sec:postcosecha}

\subsection*{Resumen de la sección}
El operador cuenta con una deshidratadora doméstica preexistente. Esta
sección queda preparada como \emph{plantilla} para registrar curvas de
secado caracterizadas experimentalmente, referencias bibliográficas de
arranque, criterios de humedad residual objetivo, presentación de
producto seco, y los hooks que conectan el lote a la base de datos
unificada del Tópico~7.

% ---------------------------------------------------------------------
\subsection{Especificaciones de la deshidratadora}
\label{sec:postcosecha:equipo}

% TODO[Operador]:
%   - Marca y modelo exactos.
%   - Potencia eléctrica [W].
%   - Número y área de bandejas [m^2 totales].
%   - Rango térmico operativo declarado por fabricante [°C].
%   - Tipo de flujo (vertical / horizontal).
%   - Sensor interno de temperatura (sí/no, precisión declarada).

\subsection{Parámetros de referencia bibliográfica}
\label{sec:postcosecha:referencias}

% TODO[Fase 1 — relleno]:
% Tabla con valores reportados en literatura, NO experimentales:
%
% | Especie | T secado [°C] | t aprox. [h] | HR residual obj. [%] | Fuente |
% |---------|---------------|--------------|----------------------|--------|
% | Oyster  | 40--50        | 6--10        | <10                  | \cite{} |
% | Lion's mane | 40--55    | 8--12        | <10                  | \cite{} |
% | Shiitake| 50--60        | 8--14        | <12                  | \cite{} |
%
% Estos valores son PUNTO DE PARTIDA. La caracterización propia se
% registra en plantilla-secado.csv y luego se vuelca a la tabla §1.6.4.
%
% Citas mínimas requeridas en referencias.bib:
%   - Chang \& Miles, "Mushrooms: Cultivation, Nutritional Value,
%     Medicinal Effect, and Environmental Impact".
%   - Stamets, "Growing Gourmet and Medicinal Mushrooms".
%   - Algún paper indexado por especie (a buscar).

\subsection{Criterio de humedad residual objetivo}
\label{sec:postcosecha:hr_residual}

% TODO[Fase 1 — relleno]:
% - Justificar rango objetivo <10--12% para conservación estable a
%   temperatura ambiente sin atmósfera modificada.
% - Indicar método de medición:
%     * Indirecto: balanza + masa antes/después.
%     * Directo (opcional, futuro): higrómetro de aw o método AOAC.
% - Validar criterio contra normativa aplicable (NMX, FDA, Codex)
%   — esta validación queda fuera de alcance Fase 1, solo se cita.

\subsection{Plantilla de caracterización (vacía)}
\label{sec:postcosecha:plantilla}

La plantilla operativa vive en
\texttt{plantilla-secado.csv} (raíz del directorio de Fase 1). Se llena
manualmente tras cada lote y, posteriormente, se importa al sistema
del Tópico~7 para construir la curva propia.

% TODO[Fase 1 — relleno]:
% Incluir aquí la misma tabla en LaTeX (vacía) para referencia impresa:
% columnas: Lote, Especie, T [°C], t [h], m_inicial [g], m_final [g],
%           %HR_residual estimada, Observaciones organolépticas.

\subsection{Flujo operativo post-cosecha}
\label{sec:postcosecha:flujo}

% TODO[Fase 1 — relleno]:
% Incluir como figura el diagrama mermaid
% diagramas/flujo-postcosecha.mmd con bifurcación fresco/seco.
% Pasos a representar:
%   - Cosecha → cepillado en seco (NO lavar con agua antes de secar)
%   - Decisión por demanda y antigüedad:
%       fresco → empaque refrigerado → distribución <72 h
%       seco   → deshidratado §1.6 → empaque hermético → distribución
%   - Registro de masa fresca cosechada en lote (Fase 7)

\subsection{Presentación de producto seco}
\label{sec:postcosecha:presentacion}

% TODO[Fase 1 — relleno]:
%   - Formatos: bolsa metalizada sellada al vacío opcional, frasco de
%     vidrio para premium.
%   - Gramajes comerciales sugeridos: 20 g, 50 g, 100 g.
%   - Etiquetado mínimo NO regulatorio (lote, fecha de deshidratado,
%     especie, recomendación de rehidratación).
%   - El detalle regulatorio (NOM-051, registros sanitarios) queda
%     fuera de alcance de la Fase 1.

\subsection{Hooks para la base de datos (Tópico 7)}
\label{sec:postcosecha:hooks}

Para que la curva propia se construya con el tiempo, la tabla
\texttt{cosechas} del schema unificado debe contemplar, por lote:

\begin{itemize}[leftmargin=*]
  \item \texttt{kg\_fresco\_cosechado}
  \item \texttt{kg\_destinado\_a\_secado}
  \item \texttt{kg\_producto\_seco\_final}
  \item \texttt{factor\_conversion} (calculado: seco / destinado)
  \item \texttt{t\_secado\_horas}, \texttt{temp\_secado\_C}
  \item \texttt{hr\_residual\_estimada\_pct}
\end{itemize}

\subsection{Criterios de aceptación}
\label{sec:postcosecha:aceptacion}

\begin{itemize}[leftmargin=*]
  \item Plantilla CSV creada y referenciada.
  \item Al menos 3 referencias bibliográficas (una por especie) en
        \texttt{referencias.bib}.
  \item Diagrama de flujo post-cosecha con bifurcación
        fresco/seco renderizable.
  \item Los campos de la tabla \texttt{cosechas} aparecen como
        requerimiento explícito para el schema del Tópico~7.
\end{itemize}
